% ============================================================
%  Paper 03: Six-Cone Partition of the Three-Sphere and Flavour Projections
%  Target: RevTeX 4-2 Compilation (SDGFT Series)
% ============================================================
\documentclass[amsmath,amssymb,aps,prd,twocolumn,superscriptaddress,nofootinbib]{revtex4-2}

\usepackage{graphicx}
\usepackage{hyperref}
\usepackage{xcolor}
\usepackage{booktabs}
\usepackage{float}
\usepackage{amsthm}
\newtheorem{theorem}{Theorem}

% ── SDGFT shared macros ──────────────────────────────────────
\usepackage{sdgft-macros}

% ── Document ─────────────────────────────────────────────────
\begin{document}

\preprint{SDGFT-2026-03}

\title{Six-Cone Partition of the Three-Sphere and Flavour Projections}

\author{David A. Besemer}
\email{david.besemer@sdgft.org}
\affiliation{Independent Researcher}

\date{\today}

\begin{abstract}
We present the mathematical proof that the three-sphere S3 is partitioned into six congruent cones with a maximum half-opening angle of 30 degrees. This geometric division represents the unique stable projection of the 24-cell onto S3. We map this partition to the flavor structure of elementary particles, showing that neutrino mixing angles are projections of this six-cone geometry.
\end{abstract}

\maketitle

\section{Introduction}
This paper details the angular partition of spatial geometry in SDGFT. The 6-cone partition serves as the template for flavor structures, matching the three generations of quarks and leptons.

\section{Theoretical Formulation}
Here we formulate the core mathematical relations governing the system:
\begin{equation}
  \Omega_{\text{cone}} = 2\pi(1 - \cos\thetamax) = 2\pi\left(1 - \frac{\sqrt{3}}{2}\right)
\end{equation}
\begin{equation}
  N_{\text{cones}} = \frac{360^\circ}{2\times 30^\circ} = 6 \quad (\text{exact spatial sectors})
\end{equation}
\section{The 30-Degree Boundary}
The boundary angle is exactly 30 degrees, corresponding to the tree-level Weinberg angle parameter. We discuss the stability of this boundary under local deformations.

\section{Conclusion}
We conclude that spatial flavor is an emergent property of the three-sphere's projective geometry under 24-cell restrictions.



\begin{thebibliography}{99}
\bibitem{Goldstein_Paper01} D. A. Besemer, ``Axiomatic Foundations of SDGFT,'' SDGFT-2026-01 (in preparation).
\bibitem{Goldstein_Paper04} D. A. Besemer, ``Effective Spectral Dimension and the Banach Fixed-Point Theorem in SDGFT,'' SDGFT-2026-04.
\bibitem{Goldstein_Code} D. A. Besemer, \texttt{sdgft} Python package, \url{https://github.com/david-besemer/sdgft} (2026).
\end{thebibliography}

\end{document}
